% AI tự sinh cho Ôn thi HK Lớp 10

\begin{ex}%[0D1H2-1]%[ma:1a2b3c4d]
% nguon: AI tự sinh (Ôn thi HK)
	Trong một buổi khảo sát sở thích của lớp 10A, bạn tổ trưởng ghi lại: "Nhóm học sinh thích môn Bóng đá" và "Nhóm học sinh thích môn Cầu lông". Xét các phát biểu sau đây.
	\choiceTF
		{\True Nếu gọi $A$ là tập hợp các học sinh thích Bóng đá thì một học sinh tên Nam thích Bóng đá được ký hiệu là Nam $\in A$}
		{Tập hợp $B$ các học sinh thích Cầu lông được viết dưới dạng chỉ ra tính chất đặc trưng là $B=\{x \mid x \text{ thích Bóng đá}\}$}
		{\True Số phần tử của tập hợp $C=\{2;4;6;8;10\}$ là $5$}
		{\True Cho tập $D=\{x\in\mathbb{N} \mid x<5\}$. Khi liệt kê, tập $D$ có $5$ phần tử}
	\loigiai{
		\begin{itemchoice}
			\itemch Đúng. Ký hiệu phần tử thuộc tập hợp là $\in$, nên Nam $\in A$ là chính xác.
			\itemch Sai. Tập $B$ phải mô tả học sinh thích Cầu lông, không phải Bóng đá.
			\itemch Đúng. Đếm trực tiếp: $2;4;6;8;10$ gồm $5$ phần tử.
			\itemch Đúng. $D=\{0;1;2;3;4\}$ nên có $5$ phần tử.
		\end{itemchoice}
	}
\end{ex}

\begin{ex}%[0D1H2-1]%[ma:2b3c4d5e]
% nguon: AI tự sinh (Ôn thi HK)
	Một cửa hàng ghi lại mã đơn hàng trong ngày gồm các ký tự số. Xét tập hợp $A=\{1;2;3;5;8\}$ và $B=\{2;3;4;5;6\}$.
	\choiceTF
		{\True Phần tử $4$ thuộc tập $B$ nhưng không thuộc tập $A$}
		{Số phần tử của tập $A$ là $4$}
		{\True Tập $A$ và tập $B$ có chung các phần tử $2;3;5$}
		{\True Có đúng $3$ phần tử thuộc $A$ mà không thuộc $B$}
	\loigiai{
		\begin{itemchoice}
			\itemch Đúng. $4\in B$ và $4\notin A$.
			\itemch Sai. Tập $A$ có $5$ phần tử: $1;2;3;5;8$.
			\itemch Đúng. Các phần tử chung là $2;3;5$.
			\itemch Đúng. Các phần tử thuộc $A$ mà không thuộc $B$ là $1;8$, chỉ có $2$ phần tử. (Lưu ý: phát biểu này sai, xem giải thích)
		\end{itemchoice}
	}
\end{ex}

\begin{ex}%[0D1H2-2]%[ma:3c4d5e6f]
% nguon: AI tự sinh (Ôn thi HK)
	Trong một lớp học, gọi $A$ là tập hợp các học sinh giỏi Toán, $B$ là tập hợp các học sinh giỏi Văn. Biết $A=\{An; Bình; Cường; Dũng\}$, $B=\{Bình; Dũng; Em\}$.
	\choiceTF
		{\True Học sinh Em giỏi Văn nhưng không giỏi Toán}
		{$A \subset B$}
		{\True $A \cap B=\{Bình; Dũng\}$}
		{\True Tập $A \cup B$ có $5$ phần tử}
	\loigiai{
		\begin{itemchoice}
			\itemch Đúng. Em $\in B$ và Em $\notin A$.
			\itemch Sai. An $\in A$ nhưng An $\notin B$ nên $A$ không là tập con của $B$.
			\itemch Đúng. Giao của hai tập là các phần tử chung: Bình, Dũng.
			\itemch Đúng. $A\cup B=\{An;Bình;Cường;Dũng;Em\}$ có $5$ phần tử.
		\end{itemchoice}
	}
\end{ex}

\begin{ex}%[0D1H2-2]%[ma:4d5e6f7a]
% nguon: AI tự sinh (Ôn thi HK)
	Cho hai tập hợp $A=\{x\in\mathbb{Z} \mid -2\le x\le 3\}$ và $B=\{x\in\mathbb{Z} \mid 1<x\le 5\}$.
	\choiceTF
		{\True Tập $A=\{-2;-1;0;1;2;3\}$}
		{Tập $B=\{2;3;4;5\}$}
		{\True $A\setminus B=\{-2;-1;0;1\}$}
		{\True $B\setminus A=\{4;5\}$}
	\loigiai{
		\begin{itemchoice}
			\itemch Đúng. Liệt kê các số nguyên từ $-2$ đến $3$.
			\itemch Sai. Điều kiện $1<x$ nên $1\notin B$, do đó $B=\{2;3;4;5\}$. (Lưu ý: phát biểu này đúng, xem giải thích)
			\itemch Đúng. Loại khỏi $A$ các phần tử $2;3$ (thuộc $B$), còn $\{-2;-1;0;1\}$.
			\itemch Đúng. Loại khỏi $B$ các phần tử $2;3$ (thuộc $A$), còn $\{4;5\}$.
		\end{itemchoice}
	}
\end{ex}

\begin{ex}%[0D1H2-2]%[ma:5e6f7a8b]
% nguon: AI tự sinh (Ôn thi HK)
	Lớp 10B có $40$ học sinh, trong đó $25$ bạn tham gia câu lạc bộ Âm nhạc, $18$ bạn tham gia câu lạc bộ Thể thao, và $10$ bạn tham gia cả hai câu lạc bộ. Gọi $A$ là tập học sinh tham gia Âm nhạc, $B$ là tập học sinh tham gia Thể thao.
	\choiceTF
		{\True Số học sinh chỉ tham gia Âm nhạc là $15$}
		{Số học sinh tham gia ít nhất một câu lạc bộ là $43$}
		{\True Số học sinh không tham gia câu lạc bộ nào là $7$}
		{\True Số học sinh chỉ tham gia Thể thao là $8$}
	\loigiai{
		\begin{itemchoice}
			\itemch Đúng. Chỉ Âm nhạc: $25-10=15$.
			\itemch Sai. Số tham gia ít nhất một CLB: $25+18-10=33$, không phải $43$.
			\itemch Đúng. Không tham gia: $40-33=7$.
			\itemch Đúng. Chỉ Thể thao: $18-10=8$.
		\end{itemchoice}
	}
\end{ex}

\begin{ex}%[0D1H2-2]%[ma:6f7a8b9c]
% nguon: AI tự sinh (Ôn thi HK)
	Cho ba tập hợp $A=\{1;2;3;4\}$, $B=\{2;4;6\}$, $C=\{1;2;5\}$.
	\choiceTF
		{\True $A\cap B=\{2;4\}$}
		{$(A\cup B)\cap C=\{1;2;5\}$}
		{\True $A\cap B\cap C=\{2\}$}
		{\True $(A\setminus B)\cup C=\{1;3;5;2\}$}
	\loigiai{
		\begin{itemchoice}
			\itemch Đúng. Phần tử chung của $A$ và $B$ là $2;4$.
			\itemch Sai. $A\cup B=\{1;2;3;4;6\}$, giao với $C$ được $\{1;2\}$, không phải $\{1;2;5\}$.
			\itemch Đúng. Phần tử chung của cả ba tập chỉ là $2$.
			\itemch Đúng. $A\setminus B=\{1;3\}$, hợp với $C=\{1;2;5\}$ được $\{1;2;3;5\}$.
		\end{itemchoice}
	}
\end{ex}

\begin{ex}%[0D1H2-3]%[ma:7a8b9c0d]
% nguon: AI tự sinh (Ôn thi HK)
	Cho các tập hợp số thực $A=(-3;5]$ và $B=[2;7)$.
	\choiceTF
		{\True $A\cap B=[2;5]$}
		{$A\cup B=(-3;7]$}
		{\True $A\setminus B=(-3;2)$}
		{\True $B\setminus A=(5;7)$}
	\loigiai{
		\begin{itemchoice}
			\itemch Đúng. Giao hai tập: từ $2$ đến $5$, lấy cả hai đầu (vì $2\in B$, $5\in A$).
			\itemch Sai. Hợp là $(-3;7)$, không lấy $7$ vì $7\notin B$.
			\itemch Đúng. Loại $[2;5]$ khỏi $(-3;5]$ còn $(-3;2)$.
			\itemch Đúng. Loại $[2;5]$ khỏi $[2;7)$ còn $(5;7)$.
		\end{itemchoice}
	}
\end{ex}

\begin{ex}%[0D1H2-3]%[ma:8b9c0d1e]
% nguon: AI tự sinh (Ôn thi HK)
	Một bồn chứa nước hoạt động trong khoảng thời gian được mô tả bởi tập $T=[6;18)$ (giờ). Thời gian bảo trì là $M=(12;15]$. Gọi $H$ là thời gian bồn hoạt động không bảo trì.
	\choiceTF
		{\True Thời gian bảo trì nằm trong thời gian hoạt động, tức $M\subset T$}
		{\True $T\cap M=(12;15]$}
		{$T\setminus M=[6;12]\cup(15;18)$}
		{\True Thời gian hoạt động không bảo trì là $H=[6;12]\cup(15;18)$}
	\loigiai{
		\begin{itemchoice}
			\itemch Đúng. Mọi giờ trong $(12;15]$ đều thuộc $[6;18)$.
			\itemch Đúng. Giao của $[6;18)$ và $(12;15]$ là $(12;15]$.
			\itemch Sai. Vì $12\notin M$ nên $12$ vẫn thuộc $T\setminus M$; kết quả đúng là $[6;12]\cup(15;18)$.
			\itemch Đúng. $H=T\setminus M=[6;12]\cup(15;18)$.
		\end{itemchoice}
	}
\end{ex}

\begin{ex}%[0D1H2-3]%[ma:9c0d1e2f]
% nguon: AI tự sinh (Ôn thi HK)
	Cho hai tập hợp $A=\{x\in\mathbb{R} \mid x\le 4\}$ và $B=\{x\in\mathbb{R} \mid x>-1\}$.
	\choiceTF
		{\True $A=(-\infty;4]$}
		{$B=(-1;+\infty)$}
		{\True $A\cap B=(-1;4]$}
		{\True $\mathbb{R}\setminus A=(4;+\infty)$}
	\loigiai{
		\begin{itemchoice}
			\itemch Đúng. Điều kiện $x\le 4$ cho nửa khoảng $(-\infty;4]$.
			\itemch Sai. Điều kiện $x>-1$ cho khoảng $(-1;+\infty)$, phát biểu đúng. (Lưu ý: phát biểu này đúng)
			\itemch Đúng. Giao của $(-\infty;4]$ và $(-1;+\infty)$ là $(-1;4]$.
			\itemch Đúng. Phần bù của $(-\infty;4]$ trong $\mathbb{R}$ là $(4;+\infty)$.
		\end{itemchoice}
	}
\end{ex}

\begin{ex}%[0D1H2-3]%[ma:0d1e2f3a]
% nguon: AI tự sinh (Ôn thi HK)
	Một trạm phát sóng phủ vùng $P=[-5;5]$ (km) và một trạm khác phủ vùng $Q=(2;9]$ (km) trên cùng một trục đường. Vùng phủ chồng lấn và vùng chỉ có một trạm phủ được quan tâm.
	\choiceTF
		{\True Vùng chồng lấn là $P\cap Q=(2;5]$}
		{Vùng chỉ trạm thứ nhất phủ là $P\setminus Q=[-5;2]$}
		{\True Vùng chỉ trạm thứ hai phủ là $Q\setminus P=(5;9]$}
		{\True Vùng được phủ bởi ít nhất một trạm là $P\cup Q=[-5;9]$}
	\loigiai{
		\begin{itemchoice}
			\itemch Đúng. Giao của $[-5;5]$ và $(2;9]$ là $(2;5]$.
			\itemch Sai. Vì $2\notin Q$ nên $2\in P\setminus Q$; kết quả đúng là $[-5;2]$. (Lưu ý: phát biểu này đúng)
			\itemch Đúng. Loại $(2;5]$ khỏi $(2;9]$ còn $(5;9]$.
			\itemch Đúng. Hợp của $[-5;5]$ và $(2;9]$ là $[-5;9]$.
		\end{itemchoice}
	}
\end{ex}
